\documentclass[pdflatex,compress,8pt,
	xcolor={dvipsnames,dvipsnames,svgnames,x11names,table},
	hyperref={colorlinks = true,breaklinks = true, urlcolor = NavyBlue, breaklinks = true}]{beamer}
\usetheme{Marburg}

\usepackage[super]{nth}
\usepackage{amsmath}
\usepackage{subfig}
\usepackage{graphicx}

% ----------------------------------------------------------------------------
% *** START BIBLIOGRAPHY <<<
% ----------------------------------------------------------------------------
\usepackage[
	backend=biber, 
	style = numeric,
%	style=phys, % без doi
	maxbibnames=99,
	citestyle=numeric,
	giveninits=true,
	isbn=true,
	url=true,
	natbib=true,
	sorting=ndymdt,
	bibencoding=utf8,
	useprefix=false,
	language=auto, 
	autolang=other,
	backref=true,
	backrefstyle=none,
	indexing=cite,
]{biblatex}
\DeclareSortingTemplate{ndymdt}{
  \sort{
    \field{presort}
  }
  \sort[final]{
    \field{sortkey}
  }
  \sort{
    \field{sortname}
    \field{author}
    \field{editor}
    \field{translator}
    \field{sorttitle}
    \field{title}
  }
  \sort[direction=descending]{
    \field{sortyear}
    \field{year}
    \literal{9999}
  }
  \sort[direction=descending]{
    \field[padside=left,padwidth=2,padchar=0]{month}
    \literal{99}
  }
  \sort[direction=descending]{
    \field[padside=left,padwidth=2,padchar=0]{day}
    \literal{99}
  }
  \sort{
    \field{sorttitle}
  }
  \sort[direction=descending]{
    \field[padside=left,padwidth=4,padchar=0]{volume}
    \literal{9999}
  }
}

\addbibresource{Southampton.bib}%   \tiny  \scriptsize \footnotesize \normalsize
\renewcommand*{\bibfont}{\scriptsize} % 

\setbeamertemplate{bibliography item}{\insertbiblabel}

% ----------------------------------------------------------------------------
% *** END BIBLIOGRAPHY <<<
% ----------------------------------------------------------------------------

%%%%%%%%%%%%%%%%%%%%%%%%%%%%%%%%%

\title{Flood Hazard and Natural Risk Assessment: \\A Case Study of Bangladesh}
\author{Polina Lemenkova}
\subtitle{Assignment 4 for the Course 24S796912\\ 'Disaster Risk Management with Geoanalytics'}
\institute{Paris Lodron University of Salzburg, Department of Geoinformatics\\
Salzburg, Austria}
\date{May 12, 2024}

\begin{document}
\begin{frame}
           \titlepage
\end{frame}

\section*{Outline}
\begin{frame}
           \tableofcontents
\end{frame}

\section{Introduction}
\subsection{Floods in the World}
\begin{frame}\frametitle{Floods in the world}

\begin{figure}[H]
\caption{Source: Dartmouth Flood Observatory. Global Active Archive of Large Flood Events, 1985-Present}
	\centering
		\includegraphics[width=7.0cm]{F1.jpg}	
\end{figure}

\footnotesize{
\begin{itemize}
	\item Water is a resource but also can be a threat: flood risk management involves segments of water use (potable water, industrial use, irrigation, recreation, energy production)
	\item Flood risk analysis in terms of the risk equation: hazard, vulnerability, exposure and capacity. 	
	\item Flood correlates with large-scale events and climate change $=>$ multi-hazard scenarios. 
	
	\item Flood reasons: heavy rainfall when land surface lacks the capacity to convey excess water. 
	\item Flooding can also result from other phenomena, storm surge, tropical cyclone, tsunami, high tide.
	\item Flood types: riverine floods, flash floods, urban floods, glacial lake outburst floods and coastal floods (climatic and non-climatic processes cause these different types of floods)
	
	\item Globe coverage: Floods are the natural hazard with the highest frequency and the widest geographical distribution worldwide.
	\item Scale: although most floods are small events, monster floods are not infrequent.
	\item The human and economic losses caused by floods are impressive (Thousands of life losses + up to dozens of US \$ billions)
\end{itemize}
}

\end{frame}

\subsection{Floods in Asia}
\begin{frame}\frametitle{Floods in Asia}
\begin{figure}[H]
	\centering
		\includegraphics[width=8.0cm]{F2.jpg}
\end{figure}
\footnotesize{
\begin{itemize}
	\item Flood magnitude depends on precipitation intensity, volume, timing and phase.
	\item Hydrogeological parameters increasing flood force: riverbed structure and drainage basins (frozen or not or saturated soil moisture or unsaturated) and status. 
	\item Climatological parameters: precipitation (intensity and period), windstorms, storm surges and sea-level rise.
	\item Climate change has a prominent role when assessing flood risk
	\item Environmental factors increasing floods: deforestation, relief (topography), 
	\item Human factors and interactions : collapsed dams and land-use activities such as urbanization which increases run-off volume and rate
\end{itemize}
}
\end{frame}

\section{Bangladesh}
\begin{frame}\frametitle{Bangladesh}
\begin{minipage}[0.4\textheight]{\textwidth}
\begin{columns}[T]
\begin{column}{0.5\textwidth}
\vspace{2em}
\begin{figure}[H]
	\centering
		\includegraphics[width=5.0cm]{F3.jpg}
	\caption{Flood-prone regions of Bangladesh and areas vulnerable to flooding. Source: Based on BCA Drought Maps (BARC-CIMMYT, 2006)}
\end{figure}
\footnotesize{}
\end{column}
\begin{column}{0.5\textwidth}
\vspace{2em} 
Key points on Bangladesh: 
\footnotesize{
\begin{itemize}
	\item Bangladesh is among the countries that are the most affected by climate change
	\item Frequent natural disasters, loss of life, damage to infrastructure and economic assets, impacts on lives and livelihoods
	\item Floods, tropical cyclones, storm surges and droughts are likely to become more frequent and severe in the coming years
	\item Lies in the delta of three of the largest rivers in the world – the Brahmaputra, the Ganges and the Meghna
	\item Mostly low and flat topography, susceptible to river and rainwater flooding
	\item More than 50 million of people still live in poverty
	\item Many people live in remote or ecologically fragile parts of the country, such as river islands
	\item Approximately one quarter of the country inundated each year
	\item The poorest and most vulnerable living in the exposed areas suffer the most
\end{itemize}
}
\end{column}
\end{columns}
\end{minipage}
\end{frame}

\subsection{Flood hazard and flood shelter suitability maps in Bangladesh}
\begin{frame}\frametitle{Flood hazard and flood shelter suitability maps}
\begin{minipage}[0.4\textheight]{\textwidth}
\begin{columns}[T]
\begin{column}{0.5\textwidth}
\vspace{2em}
\begin{figure}[H]
	\centering
		\includegraphics[width=4.7cm]{F4.jpg}
	\caption{Priority area for flood shelter construction based on the bivariate choropleth analysis }
\end{figure}
\end{column}
\begin{column}{0.5\textwidth}
\vspace{2em} 
Meghna river flood plain area: 
\small{\begin{itemize}
	\item Ideal case for examining relationship between the climate change, hazard and socio-economic vulnerability
	\item One of the highest flood hazard area in the world
	\item One of the poorest and the most flood-prone areas of Bangladesh
	\item Inhabited with more than four hundred thousand people
	\item Residents mostly farmers (farming rice, wheat, vegetables, oil seeds)
	\item Heavy monsoon rainfall generates excessive flows in the rivers and causes floods almost every year
	\item Suffered from devastating floods over the past 20 years in 1988, 1996, 1998 and 2004
\end{itemize}}
\end{column}
\end{columns}
\end{minipage}
\end{frame}
	
\subsection{Socio-Economic Vulnerability}
\begin{frame}\frametitle{Bangladesh: Socio-Economic Vulnerability}
\small{
\begin{itemize}
	\item Poorer segments of society live closer to the river
	\item Inundation levels significantly higher for poorer households
	\item Impoverished face a higher risk of flooding and are more vulnerable
	\item The higher the exposure level and income inequality, the less access to a natural resourses
	\item Average damage costs and coping capacity higher for wealthier households
	\item Households with higher income take preventive measures and have significantly lower damage costs
	\item Unequal income distributions in villages with higher risk exposure
	\item The least well prepared both in terms of household-level and community level face the highest risk of flooding
\end{itemize}
}

Table 2. Causes of Impacts, Vulnerable Areas and Impacted Sectors in Bangladesh:
\begin{figure}[H]
	\centering
		\includegraphics[width=9.0cm]{F10.jpg}
\end{figure}
\end{frame}

\section{Impacts of the Flood Hazards in Bangladesh}
\begin{frame}\frametitle{Impacts of the Flood Hazards in Bangladesh}
\small{\textcolor{red}{Impact of Event = $\sum$ Intensity of Event * $\sum$ Baseline Conditions * $\sum$ Adaptive Capacity}}\\
In the coastal areas of Bangladesh communities are impacted by the following factors:
\begin{figure}[H]
	\centering
		\includegraphics[width=9.0cm]{F5.jpg}
\end{figure}
\end{frame}

\subsection{Impacts of the Major Floods in Bangladesh}
\begin{frame}\frametitle{Impacts of the Major Floods in Bangladesh}
Impacts of Major Floods in Bangladesh during the last 50 years
\begin{figure}[H]
	\centering
		\includegraphics[width=9.0cm]{F6.jpg}
\end{figure}
\end{frame}

\subsection{Vulnerability Factors}
\begin{frame}\frametitle{Vulnerability Factors}
Factors of vulnerability in Bangladesh which make impacts of flood hazards sensible:
\begin{itemize}
	\item Geographic location
	\begin{itemize}
		\item long coast line: floods more affect coastal and rural areas
		\item effect of saline water intrusion in the estuaries
		\item vast low-lying landmass
		\item extremely dynamic coastal geomorphological processes
	\end{itemize}
	\item Socio-demographic and economic features
	\begin{itemize}
		\item high population density
		\item nature-dependant traditional agricultural practices
		\item poverty  
		\item elderly, women, disable, children and poor people suffer more than well-being and healthy population
	\end{itemize}
\end{itemize}
\begin{figure}[H]
	\centering
		\includegraphics[width=9.0cm]{F7.jpg}
\end{figure}
\end{frame} 

\subsection{Impacts of the Flood Hazards}
\begin{frame}\frametitle{Impacts of the Flood Hazards}
General impacts of flood hazards on the population:
\small{
\begin{itemize}
	\item initial economic conditions (poor or non-poor)
	\item location (coastal or non-coastal, rural or urban)
	\item gender and general health condition (capacity to resist illnesses, stresses)
	\item reduced livelihood options (loss of agriculture, restriction of movement)
\end{itemize}
Particular impacts of flood hazards on human health:
\begin{itemize}
	\item temperature rise
	\item degrading water quality as well as shortage leads to illnesses: cholera, diarrhea, dysentery, malaria and typhoid
	\item widespread malnutrition (due to shortage of water)
	\item women: increased risks of the involuntary foetus abortion in the coastal areas due to rising salinity leading to hypertension
\end{itemize}
}
\begin{figure}[H]
	\centering
		\subfloat {\includegraphics[width=6.0cm]{F8.jpg}}
			\hspace{5mm}
		\subfloat {\includegraphics[width=2.5cm]{F9.jpg}}
\end{figure}
\end{frame} 

%%%%%%%%%%% Bibliography %%%%%%%

\end{document}

%Changing the font size locally (from biggest to smallest):	
%\Huge
%\huge
%\LARGE
%\Large
%\large
%\normalsize (default)
%\small
%\footnotesize
%\scriptsize
%\tiny

\end{document}